\documentclass[11pt]{amsart}

\usepackage{amsmath, amssymb, amsthm}
\usepackage{mathtools}
\usepackage{thmtools}
\usepackage{hyperref}

% % Theorem environments
\declaretheorem[numberwithin=section]
{theorem}
\declaretheorem{lemma, proposition, corollary}[
style=plain,
sibling=theorem,
]
\declaretheorem[style=definition]
{definition, example, condition}
\declaretheorem[style=remark]{remark, note, observation}


\title{Review on \\ Intersection Homology Theory \\ \tiny{written by Mark Goresky and Robert MacPherson}}
\author{knottedtail}
\address{Chonnam University}
\email{philip94@snu.ac.kr}
\date{\today}
\keywords{}
\subjclass[2020]{} % MSC classification


\begin{document}
\maketitle

\section*{Introduction}

The main interest of this paper is a generalization of Pointcar\'e--Lefschetz intersction theory on pseudomanifolds.
The original theory can be summarized by following three propositions.
Let consider a compact oriented $n$-manifold $X$, with $i$-cycle $V$ and $j$-cycle $W$ in $X$.



\section*{Acknowledgements}
The author thanks to 

\bibliographystyle{amsplain}
% \nocite{*}
\bibliography{bibilography}

\end{document}
